% =====================================================================
% 05 — Energy-based models: motivation and the partition function.
% =====================================================================
% Planning (% comments only):
% 1. Why EBMs look appealing: global scoring of whole objects.
% 2. The partition function cost: discrete (|V|^T), continuous (R^D)^T.
% 3. Where AR and EBM differ structurally.
% 4. The Hopfield reframe: softmax attention IS an EBM update.
% =====================================================================

\section{Energy-Based Models and the Partition Function}
\label{sec:ebm_motivation}

The promise of energy-based models is to address the most central limitation that survived \cref{sec:what_holds}:
local next-token likelihood is not global correctness (\ref{cl:next_token_misalign}).
An energy-based model defines a scalar energy on whole objects and turns global compatibility into a single scoring function.
The cost is that exact normalization, exact likelihood, and exact sampling become hard.
This section states the formulation, the central obstacle, and the reframing of softmax attention as an energy-descent step.

\subsection{Energy and probability}
\label{ssec:ebm_definition}

An energy-based model parameterizes a scalar function $\Etheta : \X \to \RR$ on the object space $\X$ and induces the Boltzmann distribution
\begin{equation}
  \ptheta(x) \;=\; \frac{\exp(-\Etheta(x))}{\Ztheta},
  \qquad
  \Ztheta \;=\; \int_{\X} \exp(-\Etheta(x))\, \mathrm{d}x.
  \label{eq:ebm_definition}
\end{equation}
For finite $\X$, the integral becomes a sum.
The energy can be any neural network producing a scalar;
$\Ztheta$ is the partition function~\citep{lecun2006tutorial,song2021train}.

The appeal is that $\Etheta$ can be parameterized in any way that respects the geometry of the problem, free of the constraint that its output be a normalized probability.
A model that scores a whole sequence, a whole image, or a whole proof can be defined directly:
\begin{equation}
  \Etheta(x_{1:T}) \;=\; f_\theta(x_{1:T}),
  \label{eq:sequence_ebm}
\end{equation}
where $f_\theta$ is, for example, a transformer that returns a single scalar from a sequence input~\citep{grathwohl2019your,du2019implicit}.
Low energy denotes high compatibility;
high energy denotes incompatibility.
Two equally good completions of a prompt can have equal energy without competing for probability mass in the way they would in a locally normalized model.

\subsection{The partition function on sequences}
\label{ssec:partition_obstacle}

\Cref{eq:ebm_definition} is honest about its cost.
For a discrete sequence space $\X = \V^T$,
\begin{equation}
  \Ztheta \;=\; \sum_{x \in \V^T} \exp(-\Etheta(x)),
  \label{eq:z_discrete_seq}
\end{equation}
which contains $|\V|^T$ terms.
For a $32{,}000$-token vocabulary and a $512$-token sequence, that is a sum of $32000^{512}$ terms, well beyond enumeration.
For continuous sequence embeddings $\X = (\RR^d)^T$,
\begin{equation}
  \Ztheta \;=\; \int_{(\RR^d)^T} \exp(-\Etheta(x))\, \mathrm{d}x,
  \label{eq:z_continuous_seq}
\end{equation}
which is a high-dimensional integral whose only general guarantee is that it is hard~\citep{murphy2023advanced}.

The autoregressive factorization in \cref{eq:ar_factorization} avoided this by paying $T$ local normalizers of size $|\V|$ each instead of one global normalizer of size $|\V|^T$.
That is the structural trade.
The cost of the autoregressive bargain is the per-token mode-covering objective (\ref{cl:mode_covering}) and the local-then-global mismatch (\ref{cl:next_token_misalign}).
The cost of the energy-based bargain is \cref{eq:z_discrete_seq,eq:z_continuous_seq} or whatever approximation stands in for them.

The natural taxonomy of training methods for energy-based models reflects this:
each method makes a different choice about whether to approximate $\Ztheta$ directly, replace it with a tractable surrogate, learn only its gradient, or avoid it altogether at training time.
\Cref{sec:ebm_training} surveys these.
The natural taxonomy of inference methods reflects the same question for sampling and posterior queries;
\Cref{sec:ebm_inference} surveys those.

\subsection{The Hopfield reframing}
\label{ssec:hopfield}

A reframing that pre-dates much of the EBM discussion makes the gap between autoregressive and energy-based language modeling narrower than slogans suggest.
\citet{ramsauer2020hopfield} prove that scaled-dot-product attention, applied to query $\xi \in \RR^d$ and stored patterns $X \in \RR^{N \times d}$, is one gradient step on the continuous Hopfield energy
\begin{equation}
  E(\xi) \;=\; -\lse(\beta\, X\xi) + \tfrac{1}{2}\xi^\top \xi + c,
  \label{eq:hopfield_energy}
\end{equation}
with update
\begin{equation}
  \xi_{\mathrm{new}} \;=\; X \,\softmax(\beta\, X\xi).
  \label{eq:hopfield_update}
\end{equation}
The match is literal for a single attention head with $K = V$, with the standard $1/\sqrt{d_k}$ scaling absorbed into $\beta$;
\citet{ramsauer2020hopfield} extend it to the $K \neq V$ case by a substitution.
Capacity in this energy is exponential in pattern dimension~\citep{demircigil2017model}.

The reframe is consequential.
\begin{itemize}
  \item Softmax attention is \emph{already} an energy-descent step.
        The research question is which energy and how many steps.
  \item The local partition function \ref{cl:partition_local} is the part of $\lse$ that makes the energy admit a tractable closed-form step.
        Replacing $\lse$ with something exp-free typically removes the global competition, injective query-to-weight map, or spikiness properties that make softmax attention work~\citep{zhang2024hedgehog,han2024bridging}.
  \item The full transformer is not the gradient of any single energy.
        Multi-layer composition with layer norm, MLPs, and residuals takes the equivalence beyond literal at the layer-stack level.
        It is sound to call a single attention head an energy-descent step;
        it is metaphor to call the full transformer ``minimizing a single energy.''
\end{itemize}
This bounds the claims a project pitched against softmax can credibly make.
An EBM design that wins by ``replacing softmax with an energy'' has to specify which property of \cref{eq:hopfield_energy} the replacement preserves and which it abandons.
Recent post-mortem and synthesis papers~\citep{jelassi2024repeat,zhang2024hedgehog,han2024bridging,mongaras2025expressiveness} consistently arrive at ``softmax plus modification'' rather than ``no softmax,'' which is consistent with the energy reading.

The remaining question --- the one that the rest of the paper takes seriously --- is not ``EBM versus AR'' framed as a binary.
It is which energy, where it is evaluated, and whether $\Ztheta$ can be amortized or avoided at training and at inference.
